%%%%%%%%%%%%%%%%%%%%%%%%%%%%%%%%%%%%%%%%%%%%%%%%%%%%%%%%%%%%%%%%%%%%%%%%%%%%%%%
%% LaTeX-Vorlage für Abschlussarbeiten                                       %%
%% (TH Köln -Campus Gummersbach, Fak. 10)                                    %%
%%                                                                           %%
%% Gemäß dem Merkblatt zur Anfertigung von Projekt-, Bachelor-, Master- und  %%
%% Diplomarbeiten der Fakultät 10 von Frau Prof. Dr. Halfmann &              %%
%% Herr Prof. Dr. Rühmann (Version vom 27.01.2008)                           %%
%%                                                                           %%                                                                            
%% Bitte sprechen Sie unbedingt mit Ihrer Betreuerin bzw. Ihrem Betreuer     %%
%% bezüglich der Ausgestaltung Ihrer Arbeit!                                 %%
%%                                                                           %%
%%                                                                           %%
%% MERKKASTEN IN DIESER VORLAGE:                                             %%
%% In dieser Vorlage finden Sie Merkkasten, die Ihnen Informationen          %%
%% zu bestimmten, formalen Aspekten geben. Sprechen Sie immer auch mit       %% 
%% Ihrer Betreuerin bzw. Ihrem Betreuer dazu an.                             %%                       
%% Für die eigene Verwendung der Vorlage entfernen oder kommentieren Sie die %%
%% Merkkasten. Die betreffenden Bereiche für die Merkkasten in der Vorlage   %%
%% sind wie folgt kommentiert: <MERKKASTEN> ... </MERKKASTEN>.               %%                            %%                                                                           %%
%%                                                                           %%
%% LIZENZ:                                                                   %%
%% Diese Vorlage darf nicht kommerziell verbreitet                           %%
%% werden. Eine nicht-kommerzielle Weitergabe ist                            %% 
%% gestattet.                                                                %%
%%                                                                           %%
%% Von Ludger Schönfeld, M. Sc.,
%% 2014-2017                            %%
%%%%%%%%%%%%%%%%%%%%%%%%%%%%%%%%%%%%%%%%%%%%%%%%%%%%%%%%%%%%%%%%%%%%%%%%%%%%%%%

%%%%%%%%%%%%%%%%%%%%%%%%%%%%%%%%%%%%%%%%%%%%%
%% HEADER                                  %%
%%%%%%%%%%%%%%%%%%%%%%%%%%%%%%%%%%%%%%%%%%%%%
\documentclass[a4paper,12pt,oneside]{article}
% Optionen:
% - a4paper => DIN A4-Format
% - 12pt    => Schriftgröße (weitere  
%              grundlegende Fontgrößen: 10pt, 11pt)
% - oneside => Einseitiger Druck

%% Verwendete Pakete:
\usepackage[ngerman]{babel} % für die deutsche Sprache
\usepackage{caption} % Für schönere Bildunterschriften
\usepackage[T1]{fontenc} % Schriftkodierung (Für Sonderzeichen u.a.)
\usepackage[utf8]{inputenc} % Für die direkte Eingabe von Umlauten im Editor u.a.
\usepackage{fancyhdr} % Für Kopf- und Fußzeilen
\usepackage{lscape} % Für Querformat

%% Schriften (Beispiele)
%% Weitere LaTeX-Schriften im "LaTeX Font Catalogue"
%% unter: http://www.tug.dk/FontCatalogue/.
%% ACHTUNG: Ggf. müssen Schriften noch installiert 
%% werden!

% Serifen-Schriften:
\usepackage{lmodern} % Schriftart "Latin Modern"
%\usepackage{garamond} % Schriftart "Garamond"

%Sans Serif-Schriften:
%\usepackage[scaled]{uarial}
%\usepackage[scaled]{helvet}
%%--------------
\usepackage[normalem]{ulem} % Für das Unterstreichen von Text z.B. mit \uline{}
\usepackage[left=3cm,right=2cm,top=1.5cm,bottom=1cm,
textheight=245mm,textwidth=160mm,includeheadfoot,headsep=1cm,
footskip=1cm,headheight=14.599pt]{geometry} % Einrichtung der Seite 

\usepackage{graphicx} % Zum Laden von Graphiken
% INFO: Graphiken einbinden
%
% \includegraphics[scale=1.00]{dateiname}
%
% => Ausgabeformat: PDF-Dokument:
%    Es können die folgenden (Graphik-)formate eingebunden
%    werden: .jpg, .png, .pdf, .mps
% 
% => Ausgabeformat: DVI/PS:
%    Folgende (Graphik-)formate werden unterstützt:
%    .eps, .ps, .bmp, .pict, .pntg
\usepackage{epstopdf}

% Pakete für Tabellen
\usepackage{tabularx} % Einfache Tabellen
\usepackage{longtable} % Tabellen als Gleitobjekte (für die Aufteilung bei langen 
 %Tabellen über mehrere Seiten)
\usepackage{multirow} % Für das Verbinden von Zeilen innerhalb einer Tabelle mit
 % \multirow{anzahl}{*}{Text}

% (Zusatz-)Pakete für Formeln
\usepackage{amsmath}
\usepackage{amsthm}
\usepackage{amsfonts}

\usepackage{setspace} % Paket zum Setzen des Zeilenabstandes
% INFO: Zeilenabstand setzen:
%
% Befehle:
% - \singlespacing  => 1-zeilig (Standard)
% - \onehalfspacing => 1,5-zeilig
% - \doublespacing  => 2-zeilig 
\onehalfspacing % Zeilenabstand auf 1,5-zeilig setzen

% Farbboxen (für die Merkkästen in dieser Vorlage):
\usepackage{tcolorbox}
\tcbset{colback=white,colframe=orange,
        fonttitle=\bfseries}

\usepackage[colorlinks,pdfpagelabels,pdfstartview=FitH,
bookmarksopen=true,bookmarksnumbered=true,linkcolor=black,
plainpages=false,hypertexnames=false,citecolor=black]{hyperref} % Für Verlinkungen
% INFO: Verlinkungen mit dem hyperref-Paket:
%
% Die Angabe von URLs mit dem Befehl \url{} erlaubt einen
% gesonderten Umgang mit Weblinks. Denn die Links werden verlinkt.
% Auch erfolgt automatisch am Zeilenende ein Umbruch des Links.
% Es ist auch nicht erforderlich, Sonderzeichen in der URL manuell zu 
% entschärfen.
%
% TIPP: Sollte ein Umbuch bei einem Link nicht automatisch erfolgen, so kann
% das daran liegen, dass ein/mehrere Zeichen zusätzlich angegeben werden müssen,
% an dem der Link umbrochen werden kann.
% Dies kann mit folgendem Befehl erfolgen (Beispiel):
% \renewcommand*\UrlBreaks{\do-\do_}

% Das Paket "biblatex" für autom. 
% Literaturverzeichnisse:
%\usepackage{csquotes} % Für sprachangepasste Anführungszeichen
%\usepackage[backend=bibtex,style=alphabetic]  
%           {biblatex}
%\addbibresource{bib/literatur.bib}           

%%%%%%%%%%%%%%%%%%%%%%%%%%%%%%%%%%%%%%%%%%%%%
%% DOKUMENT                                %%
%%%%%%%%%%%%%%%%%%%%%%%%%%%%%%%%%%%%%%%%%%%%%
\begin{document}
  % Unbeschriftetes Vorblatt (Leere Seite)
  \pagestyle{empty} % Seite ohne Kopf- und Fußzeilen
  \newpage % Neue Seite
  %\newpage
\thispagestyle{empty}
\section*{}
  % Ausgelagerte LaTeX-Datei (hier: leereSeite.tex) einbinden
  
  \newpage
  
  % Deckblatt
  \pagestyle{empty}
  \begin{titlepage}
    \includegraphics[scale=1.00]{Sources/logo_TH-Koeln_CMYK_22pt}\\
    \begin{center}
      \Large
      Technische Hochschule Köln\\
      Fakultät für Informatik und Ingenieurwissenschaften\\
      \hrule\par\rule{0pt}{2cm} % Horizontale Trennlinie  mit 2 cm Abtand nach unten erzeugen
      \LARGE
      \textsc{P R A K T I K U M}\\
      \vspace{1cm} % Vertikaler Abstand von 1cm erzeugen
      \huge
      Neuronale Netze\\
      \Large
      Gruppe M.16\\
      \vspace{1.5cm}
      \large
      Vorgelegt an der TH Köln\\
      Campus Gummersbach\\
      %im Modul\\
      Mathematik 2\\ 
      \vspace{1.0cm}
      ausgearbeitet von:\\
      \begin{tabular}{c}
        \textsc{Maximilian Luca Ramacher}
      \end{tabular}
      \begin{tabular}{c c}
        \textsc{Marius Kühnast} &
        \textsc{Muratcan Garanli}\\
        \textsc{Kai Murza} &
        \textsc{Nick Struckmeyer}\\
      \end{tabular}\\
      \vspace{1.5cm}
      \begin{tabular}{ll} % Einfache Tabelle ohne Rahmen, mit 2 Spalten erzeugen
          \textbf{Erster Betreuer:} & Marc Oedingen \\
          \textbf{Zweiter Betreuer:} & Peter Wagner \\
      \end{tabular}
      \vspace{1.5cm}
      \\Gummersbach, im <Monat der Abgabe>
    \end{center}    
  \end{titlepage}
  
  \newpage
  % Abstract (ACHTUNG: Abweichung zur Reihenfolge im Merkblatt!)
  \begin{abstract}
    Platz für das deutsche Abstract...
  \end{abstract}
  
  \renewcommand{\abstractname}{Abstract}
  \begin{abstract}
    Platz für das englische Abstract...
  \end{abstract}
  %<MERKKASTEN> (für die eigene Verwendung bitte entfernen
    \vspace{1cm}
  \begin{tcolorbox}[title={Das Abstract}]
Bei einem Abstract handelt es sich um eine Art \textit{Zusammenfassung} Ihrer Arbeit. Diese kann in deutscher und/oder englischer Sprache verfasst werden. Mithilfe des Abstracts kann der Leser sich zügig orientieren, in wie fern Ihre Arbeit für ihn Relevanz besitzt.\\                                                                      Sprechen Sie unbedingt mit Ihrer Betreuerin/Ihrem Betreuer, ob Sie für Ihre Arbeit ein Abstract benötigen.\\
Ein Abstract beinhaltet folgende Aspekte \footnote{ Vgl. \cite{SW11}, S. 249}:
\begin{itemize}
 \item Ziel der Arbeit
 \item Fragestellung der Arbeit
  \item Herangezogener, theoretischer Ansatz ("Quellen")
  \item \textit{Optional:} Methodik
\end{itemize}
  \end{tcolorbox}
  %</MERKKASTEN>

%<MERKKASTEN> (für die eigene Verwendung bitte entfernen
  \vspace{1cm}
  \begin{tcolorbox}[title={Hinweise zu dieser Dokumentvorlage}]
  \begin{itemize}
   \item Es handelt sich hierbei um eine Beispiel-Vorlage für wissenschaftliche Ausarbeitungen.
Über die konkrete, formale Ausgestaltung Ihrer wissenschaftlichen Arbeit sprechen Sie unbedingt mit Ihre/m Betreuer/in.
  \item Unabhängig, ob Sie beispielsweise eine Bachelor-, Master- oder Hausarbeit schreiben müssen. Diese Vorlage kann als eine gute Basis für Ihre Arbeit dienen. Passen Sie einfach die Vorlage Ihren Anforderungen entsprechend an.
  \end{itemize}
  \end{tcolorbox}
  %</MERKKASTEN>  
  
  \newpage
  
  % Inhaltsverzeichnis
  \tableofcontents
  
  \newpage
  \pagestyle{fancy} % Kopf- und Fußzeilen aktivieren (=> Paket "fancyhdr")
 
  % Abbildungsverzeichnis  
  % INFO: Abbildung einbinden (Beispiel):
  %  \begin{figure}[h!]
  %    \centering
  %    \includegraphics[scale=1.00]{Pfad zum Bild}\\
  %    \caption{Bildunterschrift} 
  %    \label{Marke zum Referenzieren auf die Abbildung}
  %  \end{figure}
  \section*{Abbildungsverzeichnis}
  \addcontentsline{toc}{section}{Abbildungsverzeichnis} % Manuellen Eintrag im Inhaltsverzeichnis erzeugen
  \renewcommand{\listfigurename}{} % Name des Abbildungsverzeichnisses ändern
  \thispagestyle{empty}
  \listoffigures
  
  \newpage
  % Tabellenverzeichnis
  %\addcontentsline{toc}{section}{Tabellenverzeichnis}
  %\listoftables
  
  %\newpage
  
  % Hauptteil des Dokuments
  % TIPP: Jedes Kapitel oder jeder Abschnitt in eine eigene
  % LaTeX-Datei (.tex) auslagern.
  % Einbinden der ausgelagerten Dateien in diese Hauptdatei, erfolgt 
  % mittels folgendem Befehl (Beispiel):
  % beispiel.tex => \input{beispiel}
  \newpage
  % INFO: Querverweise auf Gliederungselemente, Abbildungen 
  %       & Tabellen setzen:
  %
  % Voraussetzung: Gesetzte Referenzmarke mit dem Befehl: \label{marke}
  % 
  % Referenzierung erfolgt dann mittels dem Befehl:
  % \ref{marke}
  \newpage 
  \section{Zusammenfassung und Ausblick}\label{kap_zusammfAusbl}  
   %\input{}
   TEXT FOLGT...
   
     %<MERKKASTEN> (für die eigene Verwendung bitte entfernen
    \vspace{1cm}
 \begin{tcolorbox}[title={Inhalte der \textit{Zusammenfassung und Ausblick}}]
Das Kapitel \textit{Zusammenfassung und Ausblick} enthält folgende formale Aspekte\footnote{Vgl. \cite{BBoJ},S. 6}:
\begin{itemize}
\item Kapitelweise Kurzdarstellung der Inhalte (inklusive Referenzierung auf die Kapitelnummerierung) => Nach dem Motto: \textit{Was wurde wo beschrieben?}
\item Kurzdarstellung \textit{Problem – Lösungsweg – Ergebnisse}
\item Rückkopplung auf die Einleitung: Wurde die Zielstellung der Arbeit und die Fragestellung zufriedenstellend beantwortet?
\item Kritische Bewertung (sofern nicht bereits im Hauptteil geschehen)
\item Offene Probleme
\item Richtung der zukünftigen/möglichen Arbeiten
\item Erläuterung, warum welche Aspekte in der Arbeit nicht erläutert wurden
\end{itemize}
  \end{tcolorbox}
  %</MERKKASTEN>   
   
  % Literaturverzeichnis
   % INFO: Referenzieren auf das Literaturverzeichnis:
   %
   % Befehl: \cite{refmarke}
   % 
   % "refmarke" ist die Angabe in den geschweiften Klammern bei 
   % \bibitem[]{refmarke}. 
   \newpage
    \thispagestyle{empty}
   \section{Quellenverzeichnis}
     \subsection{Literatur}
     \renewcommand{\refname}{} % Literaturverzeichnis ohne Bezeichnung
     % Variante 1: einfaches, manuelles 
     % Literaturverzeichnis
     \begin{thebibliography}{SW11} % 2. {...} => Hier die größte /breiteste Nummer (z.B. 99) oder Kurzbeleg angeben.
       \bibitem[SW11]{SW11} Stickel-Wolf, Christine; Wolf, Joachim (2011): Wissenschaftliches Lernen und Lerntechniken. Erfolgreich studieren–-gewusst wie!. Wiesbaden: Gabler. 
     \end{thebibliography} 
          
     \subsection{Internetquellen}
     \begin{thebibliography}{HR08} % 2. {...} => Hier die größte/breiteste Nummer (z.B. 99) oder Kurzbeleg angeben.
       \bibitem[BBoJ]{BBoJ}Bertelsmeier, Birgit (o. J.): Tipps zum Schrei\-b\-en ei\-n\-er Ab\-sch\-luss\-ar\-beit. Fach\-hoch\-schu\-le Köln-Campus Gummersbach, Institut für Informatik. \url{http://lwibs01.gm.fh-koeln.de/blogs/bertelsmeier/files/2008/05/abschlussarbeitsbetreuung.pdf} (29.10.2013).
        \bibitem[HR08]{HR08} Halfmann, Marion; Rühmann, Hans (2008): Merkblatt zur Anfertigung von Projekt-, Bachelor-, Master- und Diplomarbeiten der Fakultät 10. Fachhochschule Köln-Campus Gummersbach.\url{http://www.f10.fh-koeln.de/imperia/md/content/pdfs/studium/tipps/anleitungda270108.pdf} (29.10.2013).
        \bibitem[JH01]{JH01} Harrer, J. (2020). Künstliche Neuronale Netze. Gradientenabstiegsverfahren S.9
      \end{thebibliography} 
  
  % INFO: Biblatex -Ausgabe des  
  % Literaturverzeichnisses (Beispiele):   
  % - \printbibliography => Ausgabe ALLER 
  %   Einträge
  % - \printbibliography[nottype=online]
  %   => Ausgabe der Einträge, bis auf die
  %      "Online"-Einträge
  % - \printbibliography[type=online]     
  %   => Ausgabe nur der "Online"-Einträge  
   %\printbibliography

  % Anhang
  \newpage
  \setcounter{section}{0} % Nummerierung der Gliederungsebene "section" auf 0 setzen
  \renewcommand*\thesection{\Alph{section}} % Nummerierungsart für die Gliederungsebene "section" 
  % auf Großbuchstaben setzen
  \section{Anhang}\label{anhang}
    \subsection{Unterabschnitt von Anhang}\label{subsec_UabsAnhang}
    TEXT FOLGT...
  
  \newpage
  
 % Erklärung über die selbständige Abfassung der Arbeit  
 \pagestyle{empty}
 \section*{Erklärung über die selbständige\\Abfassung der Arbeit} % \section*{...}: das *-Symbol erlaubt, dass dieser
 % Gliederungspunkt nicht ins Inhaltsverzeichnis aufgenommen wird
 \addcontentsline{toc}{section}{Erklärung über die selbständige Abfassung der Arbeit}
 Ich versichere, die von mir vorgelegte Arbeit selbständig verfasst zu haben.
 Alle Stellen, die wörtlich oder sinngemäß aus veröffentlichten oder nicht veröffentlichten Arbeiten anderer entnommen sind,
 habe ich als entnommen kenntlich gemacht.\\ 
 Sämtliche Quellen und Hilfsmittel, die ich für die Arbeit benutzt habe, sind
 angegeben. Die Arbeit hat mit gleichem Inhalt bzw. in wesentlichen Teilen noch keiner anderen Prüfungsbehörde vorgelegen.\\\\
\begin{tabular}{cp{7cm}}
    & \\ 
    & \\ \hline
  \small (Ort, Datum, Unterschrift) & \normalsize \\
  \end{tabular}
  
  %<MERKKASTEN> (für die eigene Verwendung bitte entfernen
    \vspace{1cm}
 \begin{tcolorbox}[title={Hinweise zur obigen \textit{Erklärung}}]
\begin{itemize}
\item Bitte verwenden Sie nur die Erklärung, die Ihnen Ihr \textbf{Prüfungsservice} vorgibt. Ansonsten könnte es passieren, dass Ihre Abschlussarbeit nicht angenommen wird. Fragen Sie im Zweifelsfalle bei Ihrem Prüfungsservice nach.
\item Sie müssen \textbf{alle abzugebende Exemplare} Ihrer Abschlussarbeit unterzeichnen. Sonst wird die Abschlussarbeit nicht akzeptiert. 
\item Ein \textbf{Verstoß} gegen die unterzeichnete \textit{Erklärung} kann u.\,a. die Aberkennung Ihres akademischen Titels zur Folge haben.
\end{itemize}  
  \end{tcolorbox}
  %</MERKKASTEN>   
   
\newpage
% Unbeschriftetes Abschlussblatt (Leere Seite)
\thispagestyle{empty}
\newpage
\thispagestyle{empty}
\section*{}
 
    
\end{document}